\section{PDF 的重正化与匹配}
\label{sec:renorm}
%-----------------------------------------------------------------------------
本节考虑 LaMET 在计算最简单的共线 PDF 上的应用，这也是迄今为止文献中研究得最充分的部分。虽然普适性允许人们从一大类算符在大动量下的矩阵元提取共线 PDF，我们将把焦点放在与共线 PDF 极为接近的物理可观测量上，即夸克与胶子动量分布或{\it 准 PDF}。我们还将讨论坐标空间因子化途径，其中已研究了 pseudo-PDF 与流—流关联函数。

我们主要回顾使用准 PDF 在重正化与匹配方面的技术进展。匹配原则上可以在裸矩阵元层面进行，因为像 \eq{fact_0} 这样的因子化公式对裸的与重正化的动量分布都成立。裸准 PDF 中的全部 UV 发散都可以因子化进匹配系数 $C$，后者自动重正化裸格点矩阵元，随后即可取连续极限。然而，这样的匹配系数就必须在格点微扰论中计算，这既具计算挑战性又收敛缓慢~\cite{Lepage:1992xa}。更重要的是，在 UV 截断正规化下准 PDF 因 Wilson 线自能而含有线性幂次发散~\cite{Ji:2013dva,Xiong:2013bka}，使得固定阶的格点微扰论计算无法取连续极限。虽然后一问题可以通过对线性（可能还有对数）发散的重求和得到改善，但通常更可取的做法是先在格点上对准 PDF 作非微扰重正化，然后取连续极限，并在连续理论中进行微扰匹配。为此，需要彻底理解定义准 PDF 的 Wilson 线算符的可重正性。除重正化外，LaMET 的应用还依赖大动量因子化公式~\eq{fact_0} 的有效性；该式可在微扰论中被证明到所有阶，只需证明动量分布与光锥 PDF 中的共线发散相同。

我们在 \sec{renorm-nl} 中首先证明定义准 PDF 的 Wilson 线算符的乘法可重正性：先在采用 $\MS$ 方案的连续理论中工作，再把结论推广到格点理论。接着在 \sec{factorization} 中概述动量分布在微扰论所有阶的因子化定理，并给出匹配系数与 PDF 之间卷积的形式。\sec{renorm-coord} 表明因子化定理在坐标空间有等价形式，可作为从格点矩阵元提取 PDF 的另一条路径。最后，\sec{renorm-npr} 讨论准 PDF 在格点上的非微扰重正化及其与 $\MS$ PDF 的匹配。

%------------------------------------------------------------------------------
\subsection{非局域 Wilson 线算符的重正化}
\label{sec:renorm-nl}
%------------------------------------------------------------------------------
质子的动量分布由形如 Eq.~(\ref{eq:momdis1}) 的等时非局域 Wilson 线算符定义。本小节回顾这些类空非局域算符的重正化（定义 PDF 的类光非局域算符的重正化见~\cite{Collins:1981uw,Collins:2011zzd}）。我们先用辅助场方法讨论它们在 DR 中的重正化，接着讨论类似的胶子算符，再考察动量截断型 UV 正规化中的幂次发散。结论是：它们都是乘法可重正的，只与其他有限个算符发生混合。

\subsubsection{非局域夸克算符的重正化}

我们关心如下形式的算符：
\begin{equation}
	O_\Gamma(z)=\bar{\psi}\big({z\over2}\big)\Gamma W\big({z\over2},-{z\over2}\big)\psi\big(\!-\!{z\over2}\big) \ .
\end{equation}
由于 Wilson 线 $W(z_1,z_2)$ 是规范场的路径排序积分，这类算符是否乘法可重正并不显然。非类光 Wilson 圈与 Wilson 线的重正化在早期文献已有研究~\cite{Dotsenko:1979wb,Craigie:1980qs}，其乘法可重正性的全阶证明先用图形方法给出~\cite{Dotsenko:1979wb,Craigie:1980qs}，随后用规范理论的泛函形式给出~\cite{Dorn:1986dt}。同样的结论也被猜想适用于夸克双线性算符 $O_\Gamma(z)$，其重正化形式为~\cite{Musch:2010ka,Ishikawa:2016znu,Chen:2016fxx}
\begin{align} \label{eq:schemeren}
	O^B_\Gamma(z,\Lambda)=Z_{\psi,z}(\Lambda,\mu)e^{\delta m(\Lambda)|z|} O^R_\Gamma(z,\mu)\,,
\end{align}
其中``$B$''与``$R$''分别代表裸与重正化算符；$O^B_\Gamma(z,\Lambda)$ 中的场和耦合都是依赖于 UV 截断 $\Lambda$ 的裸量。$\delta m(\Lambda)$ 是 Wilson 线的``质量修正''，包含其自能的全部线性幂次发散；$Z_{\psi,z}(\Lambda,\mu)$ 则包含来自波函数与顶点重正化的全部对数发散。

早期在 $\MS$ 方案中对准 PDF 的两圈研究确实表明 $O_\Gamma(z)$ 具有乘法可重正性~\cite{Ji:2015jwa}。\eq{schemeren} 的第一个严格证明是在辅助``重夸克''场形式下给出的~\cite{Ji:2017oey,Green:2017xeu}，该形式曾被用于证明 Wilson 线的可重正性~\cite{Samuel:1978iy,Gervais:1979fv,Arefeva:1980zd,Dorn:1986dt}。这一形式通过把 QCD 拉氏量扩展以纳入辅助``重夸克''场 $Q,\bar{Q}$ 及其规范相互作用来定义：
\begin{align} \label{eq:hqet}
	{\cal L} = {\cal L}_{\rm QCD} + \bar{Q}_0in_z\cdot D_0 Q_0\,,
\end{align}
其中下标``0''表示裸量。$n_z^\mu=(0,0,0,1)$ 是类空 Wilson 线 $W(z,0)$ 的方向矢量，$D_{0}^\mu = \partial^\mu + ig_0A_{0}^\mu$，而 $Q_0$ 是 SU(3) 基础表示中的色三重态标量 Grassmann 场。注意若把 $n_z^\mu$ 换成类时矢量 $n_t^\mu=(1,0,0,0)$，\eq{hqet} 就给出 HQET 的领头阶拉氏量。

在 \eq{hqet} 定义的理论中，Wilson 线可以表达为``重夸克''场的连通两点函数
\begin{align}
	\langle Q_0(\xi)\bar{Q}_0(\eta) \rangle_Q = S_0^Q(\xi,\eta)\,,
\end{align}
其中 $\xi$、$\eta$ 是时空坐标，$\langle...\rangle_Q$ 表示对辅助场的积分。上式在相差 $in_z\cdot D_0$ 的行列式的意义下成立——该行列式是常数，可被吸收进生成泛函的归一化~\cite{Mannel:1991mc}。Green 函数 $S_0^Q(\xi,\eta)$ 满足
\begin{align}
	in_z\cdot D_0(\xi)\ S_0^Q(\xi,\eta) = \delta^{(4)}(\xi-\eta)\,,
\end{align}
选取恰当的边界条件后其解为
\begin{align}
	S_0^Q(\xi,\eta)\! =\! W(\xi^3,\eta^3)\theta(\xi^3-\eta^3)\delta(\xi^0-\eta^0)\delta^{(2)}(\vec{\xi}_\perp -\vec{\eta}_\perp)
\end{align}
这样，Wilson 线算符 $O^B_\Gamma(z)$ 可以换成两个局域复合算符之积对所有``重夸克''场位形的平均~\cite{Dorn:1986dt}：
\begin{align} \label{eq:product}
	O^B_{\Gamma}(z)& = \int d^4\xi\ \delta(\xi^3-z\big) \nn\\
	&\quad\times \langle \bar{\psi}_0\big({\xi\over2}\big) Q_0\big({\xi\over2}\big)\Gamma\bar{Q}_0\big(\!-\!{\xi\over2}\big)\psi_0\big(\!-\!{\xi\over2}\big)\rangle_Q\,,
\end{align}
UV 调节子已被省略。

于是 $O^B_{\Gamma}(z)$ 的重正化归结为两个局域``重到轻''流的重正化：
\begin{align}
	J^B = \bar{Q}_0 \psi_0\,.
\end{align}
这个辅助场理论的可重正性已用规范理论的标准泛函技术证明~\cite{Dorn:1986dt}：固定协变规范并引入鬼场之后，含辅助``重夸克''的理论具有残余 BRST 对称性，由此可导出 Ward-Takahashi 恒等式，表明 Green 函数的全部 UV 发散可用有限个局域抵消项减除。类似地，同一方法也被用于证明微扰论中 HQET 的全阶重正化~\cite{Bagan:1993zv}。

按照~\cite{Dorn:1986dt}，``重夸克''拉氏量在协变规范下可重正化为
\begin{align}\label{eq:auxl}
	{\cal L} &= {\cal L}_{\rm QCD}[g_0,\psi_0,A_0,c_0] + \bar{Q}_0in_z\cdot D_0 Q_0\nn\\
	&={\cal L}_{\rm QCD}[g,\psi,A,c] +  {\cal L}_{\rm c.t.}[g,\psi,A,c] \nn\\
	& +  Z_Q\bar{Q}\left(in_z\cdot \partial - i\delta m\right) Q - g Z_{1}^{QQg}\bar{Q}n_z\cdot A_at^aQ\,,
\end{align}
其中 ${\cal L}_{\rm c.t.}[g,\psi,A,c]$ 是 QCD 抵消项，裸场、裸耦合与重正化量的关系为
\begin{align}
	\psi_0=Z_\psi^{1\over2}\psi, \ A_0 = Z_A^{1\over2} A,\ Q_0 = Z_Q^{1\over2} Q,\ g_0=Z_gg\,.
\end{align}
重夸克—胶子顶点重正化常数 $Z_1^{QQg}$ 通过辅助场理论的 Slavnov-Taylor 恒等式与 $Z_g$ 相联系~\cite{Dorn:1986dt}：
\beq
Z_g= Z_1^{QQg} Z_A^{-{1\over2}}Z_Q^{-1}	\,.
\eeq
$i\delta m$ 可视为``重夸克''的质量修正，只不过它是虚的。对 Dirac 费米子，由于手征对称性，质量修正是对数发散且正比于裸质量的；对 HQET，重夸克的质量修正正比于 UV 截断，即线性发散，这对这里的辅助场同样成立。由于该辅助场理论的可重正性证明是在采用 DR（$d=4-2\epsilon$）的 $\MS$ 方案中进行的，所有幂次发散都消失，$\delta m$ 亦然。不过，$\delta m$ 可能因重正子模糊性含有 ${\cal O}(\Lambda_{\text{QCD}})$ 的贡献——这种模糊性已知存在于 HQET~\cite{Bigi:1994em,Beneke:1994sw}。

既然辅助场理论可重正，\eq{product} 中算符乘积的重正化就归结为两个``重到轻''流的重正化。用量子力学的标准技术~\cite{Collins:1984xc}可以递归地证明：把 $J^B$ 插入 Green 函数的整体 UV 发散在微扰论所有阶都被重正化因子 $Z_J$ 吸收：
\begin{align}
	J^B&=Z_{J}J^R = Z_{\psi}^{1/2}Z_{Q}^{1/2}Z_{V}\ J^R\,,
\end{align}
其中 $Z_{V}$ 是``重到轻''流的顶点重正化常数。HQET 中重到轻流的重正化已在微扰 QCD 中计算到三圈~\cite{Shifman:1986sm,Politzer:1988wp,Ji:1991pr,Chetyrkin:2003vi,Broadhurst:1991fz}。更近地，有论证指出``重到轻''流的反常量纲在全阶与 HQET 相同~\cite{Braun:2020ymy}；下面将讨论的``重到胶子''流也是如此，因此类空与类时 Wilson 线算符的重正化因子应当完全一致。

利用上述结果可以证明
\begin{align} \label{eq:renproof}
	O^B_\Gamma(z)\!&= Z_J^2\!\int\! d^4\xi\ \delta(\xi^3\!-\!z)\big\langle \bar{J}^R\big({\xi\over2}\big) \Gamma  J^R\big(\!-\!{\xi\over2}\big)\big\rangle_Q\nn\\
	&= Z_J^2 e^{\delta m|z|} O^R_\Gamma(z)\,,
\end{align}
其中 $\delta m$ 来自 \eq{auxl} 中 $(in_z\cdot \partial - \delta m)$ 的行列式。由此确认 \eq{schemeren} 中 $Z_{\psi,z} =  Z_J^2 $，它与 $\Gamma$ 无关。在一圈阶~\cite{Stefanis:1983ke}，
\beq \label{eq:zpsiz}
Z_{\psi,z} = 1+ \frac{\alpha_sC_F}{4\pi}{3\over\epsilon_{\mbox{\tiny UV}}}\,,
\eeq
UV 调节子 $\epsilon_{\mbox{\tiny UV}}$ 需与 DR 中的 IR 调节子 $\epsilon_{\mbox{\tiny IR}}$ 区分开。

$ O^B_\Gamma(z)$ 的乘法可重正性也通过对全阶 Feynman 图的递归分析得到证明~\cite{Ishikawa:2017faj}。除 \eq{schemeren} 外，还发现 $O^B_\Gamma(z)$ 不与胶子或其他味的夸克混合。这在辅助场形式中也容易理解，因为变味的``重到轻''流不与其他算符混合~\cite{Green:2020xco}。

最后，在格点正规化下仍可用上述技术证明 \eq{renproof}，只是此时质量修正 $\delta m$ 不为零，等于格点 UV 截断 $1/a$ 乘以耦合常数 $\alpha_s$ 的一个微扰级数。

\subsubsection{非局域胶子算符的重正化}

利用同样的``重夸克''辅助场形式，也已证明胶子准 PDF 的 Wilson 线算符是乘法可重正的~\cite{Zhang:2018diq}，图形方法的证明见~\cite{Li:2018tpe}。

按照 LaMET，胶子准 PDF 可定义为~\cite{Ji:2013dva}
\begin{align} \label{eq:quasigluon}
	\tilde g(x,P^z) = N\int {d\lambda \over 4\pi x (P^z)^2}e^{i\lambda x} \langle P| O_g(z)|P\rangle\,,
\end{align}
其中 $N$ 是归一化因子，
\begin{align}
	O^B_g(z) = g_{\perp,\mu\nu} F^{n_1\mu}_{0,a}\big({z\over2}\big) W^{ab}\big({z\over2},-{z\over2}\big) F^{n_2\nu}_{0,b}\big(\!-\!{z\over2}\big)
\end{align}
这里 $F^{n\mu}_{0,a}=n_\rho F^{\rho\mu}_{0,a}$，而 $n_1^\mu$、$n_2^\mu$ 取 $n_z^\mu$ 或 $n_t^\mu$；$a,b$ 是伴随表示的颜色指标。横向度规张量为
\begin{align}
	g_\perp^{\mu \nu} = g^{\mu \nu}  - n_t^\mu n_t^\nu / n_t^2 - n_z^\mu n_z^\nu /n_z^2\,,
\end{align}
且 $N = (n_z\cdot P/n_t\cdot P)^{(n_1+n_2)\cdot n_t}$。对格点实现，$O^B_g(z)$ 也可定义为~\cite{Dorn:1986dt,Zhang:2018diq}
\begin{align}
	O^B_g(z)\!=\! 2g_{\perp}^{\mu\nu}\mbox{tr}\Big[F^{n_1}_{0,\mu}\big({z\over2}\big)W\!\big({z\over2},\!-\!{z\over2}\big)F^{n_2}_{0,\nu}\big(\!-\!{z\over2}\big)W\!\big(\!-\!{z\over2},\!{z\over2}\big)\Big]\,,
\end{align}
其中 $F^{\mu\nu}=F^{\mu\nu}_at^a$ 而 $W$ 处于基础表示。与 \eq{product} 类似，可以把 $O^B_g(z)$ 表示为两个局域复合算符的乘积，
\begin{align} \label{eq:gluonprod}
	\tilde{O}_g^B(z)&=\int d^4\xi\ \delta(\xi^3\!-\!z)\\
	&\times g_{\perp,\mu\nu}\big\langle F^{n_1\mu}_{0,a}\big({\xi\over2}\big)Q^a_0\big({\xi\over2}\big)\bar{Q}^b_0\big(\!-\!{\xi\over2}\big)F^{n_2\nu}_{0,b}\big(\!-\!{\xi\over2}\big)\big\rangle_Q\nn\\
	&\equiv\!\int d^4\xi\ \delta(\xi^3\!-\!z) g_{\perp}^{\mu\nu}\big\langle J^{B}_{n_1\mu}\big({\xi\over2}\big)\bar{J}^{B}_{n_2\nu}\big(\!-\!{\xi\over2}\big)\big\rangle_Q\,,\nn
\end{align}
此时辅助``重''夸克场处于伴随表示，且
\beq
J^{\mu\nu}_B= F^{\mu\nu}_{0,a}Q_0^a,\ \ \ \ \bar{J}^{\mu\nu}_B= \bar{Q}_0^aF^{\mu\nu}_{0,a}\,.
\eeq

$J^{\mu\nu}_B$ 与 $\bar{J}^{\mu\nu}_B$ 的重正化比夸克情形更复杂，因为它们可与同维或更低维的其他复合算符混合。在 DR 中，BRST 对称性允许 $J^{\mu\nu}_B$
与下列算符混合~\cite{Dorn:1986dt,Zhang:2018diq}：
\begin{align}
	J^{\mu\nu}_{2B} &\!=\! \left(n_z^\nu F^{\mu n_z}_{0,a} - n_z^\mu F^{\nu n_z}_{0,a} \right) Q_0^a/n_z^2\,,\\
	J^{\mu\nu}_{3B} &\!=\! (-in_z^\mu\! A^\nu_{0,a}\! +\! in_z^\nu A^\mu_{0,a})\big[(in_z\cdot\! D_0\!-\!i\delta m)Q_0\big]^a\!/n_z^2\,.
\end{align}
其重正化矩阵为~\cite{Dorn:1986dt}
\begin{align}\label{eq:oprenormmix}
	\begin{pmatrix}
		J_{B}^{\mu\nu} \\  J_{2B}^{\mu\nu} \\ J_{3B}^{\mu\nu}
	\end{pmatrix}
	&=
	\begin{pmatrix}
		Z_{11} & Z_{12} & Z_{13}\\ 0 & Z_{22} & Z_{23} \\ 0 & 0 & Z_{33}
	\end{pmatrix}
	\begin{pmatrix}
		J_{R}^{\mu\nu}\\  J_{2R}^{\mu\nu} \\ J_{3R}^{\mu\nu}
	\end{pmatrix}
	,
\end{align}
其中 $J_{2B}^{\mu\nu}$ 规范不变，而 $J^{\mu\nu}_{3B}$ 规范依赖且正比于辅助场的运动方程（EOM）。EOM 算符的 Green 函数会产生一个 $\delta$ 函数：
\begin{align}
	(in_z\cdot D_0(\xi)-i\delta m)\langle Q_0(\xi)\bar{Q}_0(0)\rangle_Q = \delta^{(4)}(\xi)\,,
\end{align}
对辅助场积分后它只贡献接触项 $\delta(z)$。只要 $z\neq0$，这种混合在 $O^B_g(z)$ 的所有 Green 函数中都消失，因此在 $O^B_g(z)$ 的重正化中可以忽略 $J^{\mu\nu}_{B}$ 与 $J^{\mu\nu}_{3B}$ 之间的混合。当 $z=0$ 时 $O^B_g(z)$ 成为局域算符，已知会与 BRST 正合算符及 EOM 算符混合~\cite{Collins:1994ee}，其重正化可用标准方式处理。

注意当与 $n_z$ 收缩时，
\begin{align}
	J_{2B}^{n_z\mu}&\!=\! J_{B}^{n_z\mu} = F^{n_z\mu}_{0,a}Q_0^a\,,\\
	J_{3B}^{n_z\mu}&\!=\! i\big(\!-\!A^{\mu,a}_{0} + {n_z^\mu\over n_z^2} n_z\cdot A_0^a\big)\big[(in_z\cdot D_0-i\delta m)Q_0\big]_a\,,\nn
\end{align}
$J_{B}^{n_z\mu}$ 只与 EOM 算符 $J_{3B}^{n_z\mu}$ 混合；如上所述，对 $z\neq0$ 可以忽略这种混合。此外，这种简并还导致重正化矩阵元素之间的关系~\cite{Dorn:1986dt}
\begin{align}\label{eq:Zrelation}
	Z_{11}+Z_{12}=Z_{22}, \,\,\,  Z_{13}=Z_{23}\,.
\end{align}


当与 $n_t$ 收缩时，
\begin{align} \label{eq:gauxopt}
	J_{B}^{n_t\mu} &=F^{n_t\mu}_{0,a}Q_0^a\,,\nn\\
	J_{2B}^{n_t\mu} &=n_z^\mu F_{a,0}^{n_tn_z} Q_0^a/n_z^2\,,\nn\\
	J_{3B}^{n_t\mu}&= i{n_z^\mu\over n_z^2} n_t\cdot A_{0}^a \big[(in_z\cdot D_0-i\delta m)Q_0\big]_a\,.
\end{align}
可见 $J_{2B}^{n_t\mu}$ 和 $J_{3B}^{n_t\mu}$ 与 $g_{\perp}^{\mu\nu}$ 收缩后为零，因此带横向 Lorentz 指标 $\mu$ 的 $J_{B}^{n_t\mu}$ 是乘法可重正的。

综上，对 $z\neq0$ 与横向 Lorentz 指标 $\mu$，$J_{B}^{n_z\mu}$ 和 $J_{B}^{n_t\mu}$ 在坐标空间都是乘法可重正的，从而证明了胶子 Wilson 线算符 $O^B_g(z)$ 的可重正性：
\begin{align} \label{eq:renproofgluon}
	O^B_g(z)&=Z_{J}Z_{\bar{J}}\!\int\! d^4\xi\ \delta(\xi^3\!-\!z)\ g_{\perp}^{\mu\nu}\big\langle J_{n_1\mu}^{R}\big({\xi\over2}\big)\bar{J}_{n_2\nu}^{R}\big(\!-\!{\xi\over2}\big)\rangle_Q\nn\\
	&=\ e^{\delta m|z|} Z_{J}Z_{\bar{J}}\ O^R_g(z)\,,
\end{align}
其中
\begin{align}
	J_B^{n_1\mu}&=Z_{J}\ J_R^{n_1\mu}=(Z^g_Q)^{1\over 2}Z_A^{1\over 2}Z^g_V J_R^{n_1\mu}\,,\\
	\bar{J}_B^{n_2\nu}&=Z_{\bar{J}}\ J_R^{n_2\nu}=(Z^g_Q)^{1\over 2}Z_A^{1\over 2}Z^g_{\bar{V}} J_R^{n_2\mu}\,,
\end{align}
$Z^g_V$ 与 $Z^g_{\bar{V}}$ 是含一个胶子和一个``重夸克''场的顶点的重正化常数；辅助``重夸克''的波函数重正化常数 $Z^g_Q$ 因处于伴随表示而不同于夸克情形。

另外，由于量子数失配，$J_{B}^{n_z\mu}$ 和 $J_{B}^{n_t\mu}$ 不与``重到轻''夸克流混合，这意味着非局域胶子 Wilson 线算符在重正化下不与单态夸克算符混合。

\vspace{.3cm}
对极化胶子准 PDF，其定义与 \eq{quasigluon} 相同，只是胶子 Wilson 线算符变为
\begin{align}
	\Delta O^B_g(z) = \epsilon_{\perp,\mu\nu} F^{n_1\mu}_{0,a}(z) W^{ab}(z,0) F^{n_2\nu}_{0,b}(0)\,,
\end{align}
其中 $\epsilon_{\perp}^{\mu\nu}=\epsilon^{03\mu\nu}$。由于 $\epsilon_{\perp}^{\mu\nu}$ 只与横向 Lorentz 指标收缩，可以对 $O^B_g(z)$ 使用同样的证明来表明 $\Delta O^B_g(z)$ 也乘法可重正且不与单态夸克情形混合~\cite{Zhang:2018diq}。

最后，还可以证明 \eq{renproofgluon} 在格点正规化下依然成立，此时 $\delta m$ 线性发散~\cite{Zhang:2018diq}。至此完成了胶子 Wilson 线算符可重正性的证明。

\paragraph*{单圈重正化。}
现在用一个显式的单圈例子演示上述结果。要使非局域 Wilson 线算符乘法可重正，关键在于除 Wilson 线自能以外所有图中出现的线性发散必须相互抵消。为此，保持规范对称性的正规化至关重要。我们采用 DR 并保留 $d=3$ 附近的极点来辨认线性发散~\cite{Zhang:2018diq,Wang:2019tgg}。

``重到胶子''流的单圈顶点修正示于 \fig{ogmixing}。每个图的贡献为
\begin{align}\label{eq:pole_auxiliary}
	I_a^{\rho\nu}&= \frac{\alpha_s C_A}{\pi}\left[{1\over 4-d}{3\over 4} F_a^{\rho\nu}Q_a + {\rm finite\ terms}\right]\,,\nn\\
	I_b^{\rho\nu}&= \frac{\alpha_s C_A}{\pi}\Big[\frac{1}{d-4}(A_a^\nu n_z^\rho-A_a^\rho n_z^\nu) n_z\cdot \partial Q_a/n_z^2 \nn\\
	&+\frac{\pi\mu}{d-3}\big(n_z^\rho A_a^\nu-n_z^\nu A_a^\rho\big)Q_a+{\rm finite\ terms}\Big]\,,\nonumber\\
	I_c^{\rho\nu}&=\frac{\alpha_s C_A}{\pi}\Big\{\frac{1}{d-4}\Big[\frac{1}{2}\big(F_a^{\rho n_z}n_z^\nu -F_a^{\nu n_z}n_z^\rho \big) Q_a/n_z^2 \nn\\
	&+\frac{1}{4}F_a^{\rho\nu}Q_a+\frac 1 2 (A_a^\rho n_z^\nu-A_a^\nu n_z^\rho) n_z\cdot \partial Q_a/n_z^2\Big]\nonumber\\
	&-\frac{\pi\mu}{d-3}\big(n_z^\rho A_a^\nu-n_z^\nu A_a^\rho\big) Q_a+{\rm finite\ terms}\Big\}\,.
\end{align}

\begin{figure}[htb]
	\centering
	\begin{subfigure}{.3\linewidth}
		\centering
		\includegraphics[width=\linewidth]{images/pdf-Qg_fig07a.pdf}
		\caption{}
		\label{fig:Qgv1}
	\end{subfigure}
	\begin{subfigure}{.3\linewidth}
		\centering
		\includegraphics[width=\linewidth]{images/pdf-Qg_fig07b.pdf}
		\caption{}
		\label{fig:Qgv2}
	\end{subfigure}
	\begin{subfigure}{.3\linewidth}
		\centering
		\includegraphics[width=\linewidth]{images/pdf-Qg_fig07c.pdf}
		\caption{}
		\label{fig:Qgv3}
	\end{subfigure}
	\caption{``重到胶子''流的单圈顶点修正。}
	\label{fig:ogmixing}
\end{figure}

\fig{Qgv2} 与 \fig{Qgv3} 都含有明显的线性发散（即 $\mu/(d-3)$ 项），但二者相互抵消。这保证了顶点修正的整体 UV 发散是对数型的，从而``重到胶子''流的重正化是乘法的。合并 \eq{pole_auxiliary} 的单圈结果与波函数重正化，可得
\begin{align}
	Z_{11} &= 1+ \frac{\alpha_s C_A}{4\pi}{1\over \epsilon_{\mbox{\tiny UV}}}\,,\qquad\qquad Z_{12} = 1- \frac{\alpha_s C_A}{4\pi}{1\over \epsilon_{\mbox{\tiny UV}}}\,,\nn\\
	Z_{13} &=Z_{23} = 1- \frac{\alpha_s C_A}{4\pi}{1\over \epsilon_{\mbox{\tiny UV}}}\,,\quad Z_{22} = 0\,,
\end{align}
其中 QCD 的 $C_A=N_c=3$。
如果忽略到 EOM 算符的混合，
\begin{align}
	Z^{J^{n_z\nu}}_V &= Z^{J^{\nu n_z}}_V = 0 \,,\nn\\
	Z^{J^{n_ti}}_V &= Z^{J^{in_t}}_V = Z^{J^{ij}}_V = Z^{J^{ji}}_V = 1+ \frac{\alpha_s C_A}{4\pi}  {1\over \epsilon_{\mbox{\tiny UV}}}\,,
\end{align}
$i,j=1,2$。于是单圈流重正化常数为
\begin{align}
	Z_{J^{n_z\nu}} &= Z_{J^{\nu n_z}} =1 +{\alpha_s\over 4\pi}\left({1\over 6}C_A - {4\over 3}n_fT_F\right) {1\over\epsilon_{\mbox{\tiny UV}}} \,,\nn\\
	Z_{J^{n_ti}} &= Z_{J^{in_t}} = Z_{J^{ij}} =Z_{J^{ji}} \nn\\
	&=1 + {\alpha_s\over 4\pi}\left({7\over 6}C_A - {4\over 3}n_fT_F\right) {1\over\epsilon_{\mbox{\tiny UV}}}\,,
\end{align}
$T_F=1/2$，$n_f$ 为活跃夸克味数。两圈结果见~\cite{Braun:2020ymy}。

可以看到，对 $\mu,\nu=0,1,2$，``重到胶子''流的反常量纲相同，这源于围绕 $z$ 轴的 $SO(2,1)$（欧氏空间中为 $SO(3)$）对称性。

%------------------------------------------------------------------------------
\subsection{准 PDF 的因子化}
\label{sec:factorization}
%------------------------------------------------------------------------------


LaMET 应用于共线部分子物理的关键是把准 PDF 与光锥 PDF 相联系的因子化公式~\cite{Ji:2013dva}。这里我们利用匹配系数的微扰性质，把 $\MS$ 方案中的因子化形式写成与任意给定 $x$ 处 PDF 直接 EFT 计算相一致的方式~\cite{Izubuchi:2018srq,Wang:2019tgg}：
\begin{align}
	q_i(x,  \mu)\! &=\! \int_{-\infty}^{\infty} \frac{dy}{|y|} \ \bigg[\!\sum_j \tilde C_{q_iq_j}\!\left(\frac{x}{y}, \frac{\mu}{yP^z}\right) \tilde{q}_j(y,P^z,\mu) \label{eq:factorization1} \\
	&\qquad + \tilde C_{qg}\left(\frac{x}{y}, \frac{\mu}{yP^z}\right) \tilde{g}(x,P^z,\mu)\bigg] + \cdots
	\ ,\nn\\
	g(x, \mu)\! &=\! \int_{-\infty}^{\infty} \frac{dy}{|y|} \ \bigg[\!\sum_j \tilde C_{gq}\left(\frac{x}{y}, \frac{\mu}{yP^z}\right) \tilde{q}_j(y,P^z,\mu) \label{eq:factorization2} \\
	&\qquad + \tilde C_{gg}\left(\frac{x}{y}, \frac{\mu}{yP^z}\right) \tilde{g}(y,P^z,\mu)\bigg] + \cdots\,,\nn
\end{align}
其中 $i,j$ 取遍夸克与反夸克的味。``$\cdots$''项包括目标质量修正（其解析形式已在所有 $M^2/(P^z)^2$ 阶导出~\cite{Chen:2016utp}）以及 ${\cal O}\big(\Lambda_{\text{QCD}}^2/ (xP^z)^2,\Lambda_{\text{QCD}}^2/ ((1-x)P^z)^2\big)$ 阶的高扭度贡献（见 \eq{fact_0}）。右端所有的 $P^z$ 依赖都会消去，正如重正化标度那样。

如 \sec{lamet} 所解释，上述因子化有坚实的物理基础：准 PDF 与光锥 PDF 的差别只是 $P^z\to \infty$ 与 $\Lambda_{\rm UV}\to \infty$ 的极限次序，二者的 IR 物理必然相同。夸克准 PDF 在微扰论中的全阶因子化证明最早由图形方法给出~\cite{Ma:2014jla}；该公式也由非局域 Wilson 线算符的算符乘积展开（OPE）导出过~\cite{Izubuchi:2018srq,Ma:2017pxb,Wang:2019tgg}。此处我们概述类似于~\cite{Ma:2014jla}的图形证明，说明准 PDF 的共线发散确实可因子化且与光锥 PDF 的共线发散相等。由于共线发散是微扰论中的概念，我们将用以类光动量 $P^\mu=(P^z,0,0,P^z)$ 的无质量外部夸克态来展示因子化。虽然证明只针对微扰的自由夸克态，但因子化公式被广泛相信在非微扰层面同样成立。我们用 DR 同时调节 UV 与共线发散，并且只考虑裸量，因为 UV 重正化不改变领先的共线发散。

在分析之前应指出：如 \sec{lamet-rg} 所述，实修正与虚修正对准 PDF 的软发散完全相消，故只需关注共线发散。为直观理解共线发散的结构，我们从 Feynman 规范下的单圈图 \fig{qoneloop1} 出发。积分为
\begin{align}
	\int \frac{d^{4-2\epsilon}k}{(2\pi)^{4-2\epsilon}}\frac{{\rm tr}\big[\slashed{P}\slashed{k}\gamma^z\slashed{k}\big]\delta(k^z-yP^z)}{(k^2+i0)^2((P-k)^2+i0)}\,.
\end{align}
内线夸克动量为 $k^\mu = (k^+,k^-,\vec{k}_\perp)$，胶子动量为 $P-k$。当 $k^-$ 与 $k_\perp=|\vec{k}_\perp|$ 很小时，内线夸克与胶子变得与外部夸克共线，即 $k^{\mu}\sim (k^+,0,0_\perp)$ 且 $(P-k)^{\mu}\sim (P^+-k^+,0,0_\perp)$。此时夸克与胶子传播子的分母 $(k^2)^2$ 与 $(P-k)^2$ 都趋于零，导致共线发散。反过来，若 $k^2=(P-k)^2=0$，则条件要求 $k^2=k\cdot P=P^2=0$，故 $k$ 必与 $P$ 共线。对小 $k^-$ 与 $k_\perp$，$\delta$ 函数由 $k^z=(k^+-k^-)/\sqrt{2}$ 中的 $k^+$ 项主导，约化为 $\sqrt{2}\delta(k^+-yP^+)$。这正是限制光锥 PDF 中 $k^+=yP^+$ 的顶点（差一个因子 $\sqrt{2}$）。再者，对共线的 $k$ 与 $(P-k)$，分子中的旋量迹由 $\gamma^z=(\gamma^+-\gamma^-)/\sqrt{2}$ 的 $\gamma^+$ 部分主导：${\rm tr}\big[\slashed{P}\slashed{k}\gamma^z\slashed{k}\big] \sim {\rm tr}\big[\slashed{P}\slashed{k}\gamma^+\slashed{k}\big]/\sqrt{2}$。于是，在共线区 $k^{\mu}\sim (k^+,0,0,0)$，上述积分约化为光锥 PDF 的积分：
\begin{align}
	\int_{\rm c} \frac{d^{4-2\epsilon}k}{(2\pi)^{4-2\epsilon}}\frac{{\rm tr}\big[\slashed{P}\slashed{k}\gamma^+\slashed{k}\big]\delta(k^+-yP^+) }{(k^2+i0)^2((P-k)^2+i0)}\,,
\end{align}
下标``c''表示共线区。

上述图像自然出现在高度推进的强子态中，那里夸克近似在壳。因此，正如 \sec{lamet-univ} 所解释的，{\it 尽管算符本身不含光锥信息，大动量的外部强子态仍能生成与光锥 PDF 等价的共线发散。}把准 PDF 的完整积分减去光锥 PDF 的对应积分，对数的共线发散便相消，剩下的差别是微扰的，可被吸收进匹配核。

类似地，对 Fig.~\ref{fig:qoneloop2a} 的顶点图，圈积分正比于
\begin{align}
	\int \frac{d^{4-2\epsilon}k}{(2\pi)^{4-2\epsilon}}\frac{1}{P^z-k^z}\frac{{\rm tr}\big[ \slashed{P} \gamma^z \slashed{k} \gamma^z\big] \delta(k^z-yP^z)}{(k^2+i0)((P-k)^2+i0)} \,.
\end{align}
整个积分在共线区约化为
\begin{align}
	\int_{\rm c} \frac{d^{4-2\epsilon}k}{(2\pi)^{4-2\epsilon}}\frac{1}{P^+-k^+}\frac{{\rm tr} \big[\slashed{P} \gamma^+ \slashed{k} \gamma^+\big]\delta(k^+-yP^+)}{(k^2+i0)((P-k)^2+i0)}  \ ,
\end{align}
这正是光锥 PDF 的相应积分。该图的一个关键特征是：虽然规范链探测的是胶子场的 $z$ 分量 $A^z =\left(A^+-A^-\right)/\sqrt{2}$，但对领先共线发散有贡献的只有 $A^+$ 分量（纵向极化）。往规范链上再接一条共线胶子会通过链传播子带来 ${\cal O}(1/P^z)$ 的幂次压低，但从快速运动的色荷辐射出的共线胶子的 $A^+$ 分量获得 Lorentz 推进因子 $\gamma$ 的增强，恰好补偿了压低。

上述结果可以推广到所有阶。与单圈图类似，在共线发散的领先区存在任意多条纵向极化的 $A^+$ 胶子，它们是从快速运动态动力学发射出来的，而不像标准共线 PDF 那样靠类光规范链手工放入。$A^+$ 胶子的存在显然增加了展示准 PDF 与光锥 PDF 共线发散等价性的复杂度。为简化起见，以下改在光锥规范 $A^+=0$ 下工作以消除所有 $A^+$ 胶子。这样顶点图不再对领先共线发散有贡献，使其结构大大简化。

在一般图形中，我们把准 PDF 的潜在领先区分解成 Fig.~\ref{fig:ladder} 所示的梯形结构。包含定义准 PDF 的非局域算符的上部两粒子不可约（2PI）核记为 $H$；梯中的 2PI 核记为 $K$。$K$ 含上部两条外夸克线但不含下部的。当有 $n$ 个 2PI 核时，流出梯子的动量自下而上标记为 $k_1$ 到 $k_n$。把 $H$ 与 $K$ 写成旋量空间与动量空间中的矩阵：$H=H_{\alpha'\beta'}(yP^z;k)$，其中 $k$ 为流入 $H$ 的动量；$K=K_{\alpha\beta;\alpha'\beta'}(k,k')$，其中 $k$、$k'$ 分别为上下外腿的动量。$\alpha\beta$ 与 $\alpha'\beta'$ 分别是上下两条外腿的旋量指标。
\begin{figure}
	\includegraphics[width=0.8\columnwidth]{images/pdf-ladder_fig08.pdf}
	\caption{\label{fig:ladder} 准 PDF 的梯分解（左）。上部 2PI 核 $H$ 包含定义准 PDF 的算符，图底部的外腿是外部大 $P^z$ 态。右侧显示核 $H$ 与 $K$。}
\end{figure}
循着~\cite{Curci:1980uw,Collins:2011zzd}的方法，我们发现：
\begin{enumerate}
	\item 在光锥规范下，上部 $H$ 中没有共线发散。
	\item 若 $k_1,\cdots,k_n$ 都不共线，则没有领先共线发散。更一般地，对第 $i$ 个 2PI 核，只要 $k_{i-1}$ 与 $k_{i}$ 中有一个不共线，核内部的子积分就有限，不对领先共线发散作贡献。
	\item 若 $k_i$ 不共线，则第 $i$ 个梯以上部分的图形没有共线发散。
\end{enumerate}
因此，共线发散产生于动量区域 $R_i$：其中 $k_1$ 到 $k_i$ 共线而 $k_{i+1}$ 到 $k_n$ 不共线。我们可以构造抵消项来减除每个区域 $R_i$ 中的共线发散。为此，如同单圈例子那样，在收敛的上部 $HK^{n-i}$ 中只保留 $k_i$ 的 $+$ 分量，即在上部取 $k_i\rightarrow (k^+_i,0,0_\perp)$。这显然不会改变共线发散。还要注意 $[HK^{n-i}]_{\alpha\beta}=H_{\alpha'\beta'}K^{n-i}_{\alpha'\beta';\alpha\beta}$ 应理解为指标为 $\alpha\beta$ 的 $4\times 4$ Dirac 矩阵，而下部是 $[K^i \slashed P]_{\alpha\beta}=K^i_{\alpha\beta;\alpha'\beta'}\slashed P_{\alpha'\beta'}$。在共线发散的领先区，$HK^{n-i}$ 与 $K^i\slashed P$ 分别正比于 $\gamma^+$ 与 $\gamma^-$。
于是，为获得领先共线发散，可以分别用 $\gamma^-/2$ 和 $\gamma^+/2$ 收缩上下两部分，把二者的旋量迹解开；二者之间唯一的沟通是 $k^+$ 积分。共线发散包含在下部
\begin{align}
	q^{i}(x,\epsilon_{\mbox{\tiny IR}})\!=\!\!\int \frac{dk^-\!d^{d-2}\!k_\perp}{2(2\pi)^{d}}{\rm tr} \big[\gamma^+\! K^i(xP^+\!,k^-\!,k_\perp;P)\slashed{P}\big],
\end{align}
其中 $d=4-2\epsilon$，$k^+=xP^+$。区域 $R_i$ 的减除可以有效地写成一个卷积
\begin{align}
	\int \frac{dx}{x} \hat C^{n-i}(y,x,P^z)q^i(x,\epsilon_{\mbox{\tiny IR}}) \ ,
\end{align}
其中
\begin{align}
	\hat C^{n-i}(y,x,\!P^z)=\frac{1}{2}{\rm tr}\left[HK^{n-i}(yP^z;xP^+,\!0,0_\perp)(xP^+)\gamma^-\right]
\end{align}
是朴素匹配核，$y$ 依赖来自 $H$ 中的算符。然而朴素匹配核仍带有需要减除的共线子发散。这可以通过递归定义的减除匹配核 $C^{n-i}(y,x)$ 实现，方式类似于 UV 重正化的 BPHZ 关系~\cite{Collins:2011zzd}。对 $n$ 与 $i$ 求和后，递推关系给出
\begin{align}\label{eq:subtraction}
	\tilde q(y,P^z,\epsilon_{\mbox{\tiny IR}})&=\sum_{n=0}^{\infty}\sum_{i=0}^n \int \frac{dx}{x} C^{n-i}(y,x,P^z)q^i(x,\epsilon_{\mbox{\tiny IR}})\nn\\
	&=\int \frac{dx}{x} C(y,x,P^z)q(x,\epsilon_{\mbox{\tiny IR}}) \ ,
\end{align}
其中 $\tilde q(y,P^z,\epsilon_{\mbox{\tiny IR}})$ 是准 PDF，$C(y,x,P^z)=\sum_{n=0}^{\infty} C^{n}(y,x,P^z)$ 是全阶匹配核，$q(x,\epsilon_{\mbox{\tiny IR}})=\sum_{i=0} q^i(x,\epsilon_{\mbox{\tiny IR}})$。基于 $q^i(x,\epsilon_{\mbox{\tiny IR}})$ 的定义，显然 $q^i$ 就是带 $i$ 个 2PI 核的光锥 PDF，而 $q$ 是具有自然支撑 $0<x<1$ 的完整光锥 PDF。{\it 光锥 PDF $q(x)$ 与定义准 PDF 的算符无关}，因为它只对共线发散的显式形式敏感。\eq{subtraction} 右端包含准 PDF $\tilde q$ 的全部共线发散，于是裸量之间的匹配关系得以确立。对重正化后的量可以写出类似的匹配，此时重正化只影响匹配核 $C(y,x,P^z)$。注意，导致 \eq{fact_0} 的显式解 $C^{n-i}(y,x,P^z)$ 可以基于与~\cite{Collins:2011zzd}中类似的减除算符给出。另外，\eq{subtraction} 可按 $\alpha_s$ 逐阶反解，从而证明了 \eq{fact_0}，并可推广到 \eqs{factorization1}{factorization2}。

现在给出 $\MS$ 方案下单圈的匹配系数。动量 $p^\mu=(p^z,0,0,p^z)$ 的自由无质量夸克态中 $\MS$ 准 PDF 与光锥 PDF 的单圈展开为
\begin{align}
	\tilde{q}(y,\mu/p^z,\epsilon_{\mbox{\tiny IR}}) &= \tilde{q}^{(0)}(y) + {\alpha_sC_F\over 2\pi}\tilde{q}^{(1)}(y,\mu/p^z,\epsilon_{\mbox{\tiny IR}})\,,\\
	q(x,\epsilon_{\mbox{\tiny IR}}) &= q^{(0)}(x) + {\alpha_sC_F\over 2\pi}q^{(1)}(x,\epsilon_{\mbox{\tiny IR}})\,.
\end{align}
树图水平上 $\tilde{q}^{(0)}(y) = q^{(0)}(y) = \delta(1-y)$。一圈时，$\MS$ 准 PDF 及其抵消项为~\cite{Izubuchi:2018srq}
\begin{align}\label{eq:qPDFren}
	&\tilde{q}^{(1)}(y,\mu/p^z,\epsilon_{\mbox{\tiny IR}}) \nn\\
	=&\left\{
	\begin{array}{ll}
		\Big({1+y^2\over 1-y}\ln {y\over y-1} + 1 + {3\over 2y}\Big)^{[1,\infty]}_{+(1)}- {3\over 2y} & y>1\\
		\Big({1+y^2\over 1-y}\big[- {1\over \epsilon_{\mbox{\tiny IR}}} - \ln{\mu^2\over 4(p^z)^2} + \ln\big(y(1-y)\big)\big] &\\
		\  - {y(1+y)\over 1-y}+ 2\sigma (1-y)\Big)^{[0,1]}_{+(1)}  &0<y<1\\
		\Big(-{1+y^2\over 1-y}\ln {-y\over 1-y} - 1 + {3\over 2(1-y)}\Big)^{[-\infty,0]}_{+(1)}&\\
		\ - {3\over 2(1-y)} & y<0
	\end{array}\right.\nn\\
	&+ \delta(1-y)\left[{3\over2}\ln{\mu^2\over 4(p^z)^2} + {5 + 2\sigma\over2}\right]\,,
\end{align}
\begin{align}
	\delta \tilde{q}^{(1)}(y,\mu/p^z,\epsilon_{\mbox{\tiny UV}})
	= {3\over 2\epsilon_{\mbox{\tiny UV}}}
	\delta(1-y)\,,
\end{align}
其中 $\epsilon_{\mbox{\tiny IR}}$ 调节共线发散；$\sigma=0$ 对应 $\Gamma=\gamma^t$，$\sigma=1$ 对应 $\Gamma=\gamma^z$。支撑在区域 $D$ 内、位于 $y=y_0$ 的加号函数定义为
\begin{align}
	\int_D\! dy \big[ g(y)\big]_{+(y_0)}^D h(y)\! =\! \int_D\! dy\ g(y) \left[ h(y)\! -\! h(y_0)\right]
\end{align}
（$g(y)$、$h(y)$ 任意）。注意准 PDF 的 $\MS$ 重正化实际上需要对矢量流守恒作细致处理~\cite{Izubuchi:2018srq}；这里只给出足以支撑我们讨论的形式，与~\cite{Izubuchi:2018srq}的差异仅在 $y=\pm\infty$ 处的 $\delta$ 函数，与~\cite{Alexandrou:2019lfo}的处理亦略有不同。

另一方面，
\beq\label{eq:lcpdf}
q^{(1)}(x,\epsilon_{\mbox{\tiny IR}}) ={\alpha_sC_F\over 2\pi} {(-1)\over \epsilon_{\mbox{\tiny IR}}} \left({1+x^2\over 1-x}\right)^{[0,1]}_{+(1)}\,,
\eeq
如预期的那样局限于物理区域。

比较 \eqs{qPDFren}{lcpdf} 中的准 PDF 与光锥 PDF，发现二者具有相同的共线发散；换言之，它们共享同样的 IR 物理，从而在一圈阶验证了因子化公式。取 $p^z=xP^z$ 并把一圈结果代入 \eq{fact_0}，即可提取只依赖微扰标度 $\mu$ 和 $P^z$ 的强子矩阵元匹配系数：
\begin{align}\label{eq:msbarmatching}
	C^{\MS}\left(y, {\mu\over xP^z}\right) &= \delta\left(1-y\right) + {\alpha_sC_F\over 2\pi}\left[\tilde q^{(1)}\left(y,{\mu\over xP^z}, \epsilon_{\mbox{\tiny IR}}\right)\right.\nn\\
	&\qquad\left.-q^{(1)}(y,\epsilon_{\mbox{\tiny IR}})\right]\,.
\end{align}
横动量截断方案与 $\MS$ 方案下 \eq{fact_0} 的完整一圈匹配系数见 \cite{Wang:2017qyg,Wang:2017eel,Wang:2019tgg}；两圈结果最近在~\cite{Chen:2020arf,Chen:2020iqi,Chen:2020ody,Li:2020xml}中得到。

%------------------------------------------------------------------------------
\subsection{双线性算符的坐标空间因子化}
\label{sec:renorm-coord}
%------------------------------------------------------------------------------
尽管 LaMET 对 PDF 的应用关注 $P^z\to\infty$ 极限下动量密度的展开，格点 QCD 计算实际上是从坐标空间关联出发的，例如
\beq
\tilde h(z, P^z) = {1\over N_\Gamma}\langle P^z| O_\Gamma(z) |P^z\rangle\,,
\eeq
先在所有 $z$ 上计算，再在固定 $P^z$ 下对 $\lambda=zP^z$ 作 Fourier 变换。这里归一化因子 $N_\Gamma=2P^z$（$\Gamma=\gamma^z$）、$N_\Gamma=2P^t$（$\Gamma=\gamma^t$）。$\tilde h(z,P^z)$ 是两个独立变量 $z$ 与 $P^z$ 的函数；在 LaMET 分析中相关的组合是准 LF 距离 $\lambda$（见 \fig{Boost}）与 $P^z$，因此 $\tilde h(\lambda, P^z)$ 将被称为准 LF 关联，以区别于下面的 LF 关联 $h(\lambda,\mu)$。

文献~\cite{Braun:1994jq} 的坐标空间因子化被建议作为从 $\tilde h(z, P^z)$ 提取 PDF 的另一条路径~\cite{Radyushkin:2017cyf,Orginos:2017kos,Radyushkin:2019mye}，它与 OPE 密切相关。不必以 $\lambda$ 和 $P^z$ 为变量，也可以把 $\tilde h$ 视为 $\lambda$ 与 $z^2$ 的函数，即 $\tilde h(\lambda, z^2)$。$\tilde h(\lambda, z^2)$ 对 $\lambda$ 的 Fourier 变换不再是质子在固定动量下的动量分布，而是所谓的赝分布（pseudo-distribution）~\cite{Radyushkin:2017cyf}。在小 $|z|\ll \Lambda^{-1}_{\rm QCD}$ 处，$\tilde h(\lambda, z^2)$ 可因子化为光锥关联~\cite{Radyushkin:2017lvu,Izubuchi:2018srq}：
\begin{align}
	\label{eq:io-Q-fact}
	\tilde{h}(\lambda, z^2\mu^2)
	=\int_{-1}^1 d\alpha\ \mathcal{C}(\alpha, z^2\mu^2)\ h(\alpha\lambda,\mu) +...\,,
\end{align}
其中 $...$ 是 $z^2\Lambda^2_{\rm QCD}$ 的幂次修正，匹配系数 $\mathcal{C}$ 与 Eq.~(\ref{eq:fact_0}) 中的 $C$ 通过
\begin{align} \label{eq:ftc-C}
	C\left(\eta,{\mu\over xP^z}\right)
	\!= \!\int {d\lambda\over 2\pi}\ e^{i\eta\lambda}\int_{-1}^1\!\! d\alpha\: e^{-i\lambda\alpha}\: \mathcal{C}\left(\alpha,{\mu^2 \lambda^2\over (xP^z)^2}\right)\,.
\end{align}
相联系。

为说明 \eq{io-Q-fact} 的因子化与 OPE 的联系，以非单态夸克为例~\cite{Izubuchi:2018srq,Wang:2019tgg}。在 $\MS$ 方案中，重正化的 $O_{\gamma^{\mu_0}}(z,\mu)$ 在 $z^2\to 0$ 时可展开为局域规范不变的扭度二算符级数：
\begin{align}\label{eq:ope-tq}
	O_{\gamma^{\mu_0}}(z,\mu) &=\sum_{n=0}^\infty \Big[C_n ({\mu}^2 z^2)\frac{(iz)^n}{n!} (n_z)_{\mu_1}\cdots (n_z)_{\mu_n} \nn\\
	&\qquad \times O^{\mu_0\mu_1\cdots\mu_n}(\mu) +  \text{higher-twist}\Big]\,,
\end{align}
其中 $\mu_0\!=\!0$ 或 $3$，$C_n\!=\!1+{\cal O}(\alpha_s)$ 是 Wilson 系数，$O^{\mu_0\mu_1\cdots\mu_n}(\mu)$ 即 \eq{t2opq} 的扭度二算符。

利用 \eq{t2mom} 的强子矩阵元及其与光锥 PDF 的关系 \eq{t2sum}，可写出 $O_{\gamma^{\mu_0}}(z,\mu) $ 强子矩阵元的小 $|z|$ 展开~\cite{Izubuchi:2018srq}：
\begin{align}\label{eq:ope}
	\tilde{h}(\lambda,z^2\mu^2)&=\langle P| O_{\gamma^{\mu_0}}(z,\mu) |P\rangle/(2P^{\mu_0}) \nn\\
	&= \sum_{n=0}^\infty\ C_n(z^2\mu^2){(-i\lambda)^n\over n!}\left[1+{\cal O}\big({M^2\over (P^z)^2}\big)\right]\nn\\
	&\quad\times \int_{-1}^1 dx\ x^n q(x,\mu) + {\cal O}\left(z^2\Lambda_{\text{QCD}}^2\right)\,,
\end{align}
其中 ${\cal O}\left(M^2/(P^z)^2\right)$ 项来自运动学迹的贡献，${\cal O}\left(z^2\Lambda_{\text{QCD}}^2\right)$ 项来自高扭度。Wilson 系数 $C_n(z^2\mu^2)$ 已在一圈~\cite{Izubuchi:2018srq}与两圈~\cite{Li:2020xml}阶计算。比较上式与 Eq.~(\ref{eq:io-Q-fact})，可以认定
\begin{align}  \label{eq:fac-C}
	\mathcal{C} (\alpha, {\mu}^2 z^2) \equiv \int\! \frac{d \lambda}{2 \pi} \,
	e^{i  \lambda \alpha}  \sum_n C_n ({\mu}^2 z^2)  \frac{(-i \lambda)^n}{n!}
	\,.
\end{align}
由于 $z^2$ 在 $\mathcal{C} (\alpha, {\mu}^2 z^2)$ 中固定，\eq{fac-C} 中的积分实际上是对 $P^z$ 从 $-\infty$ 到 $+\infty$ 进行的。$\mathcal{C}(\alpha,z^2\mu^2)$ 的支撑为 $-1 \le \alpha \le 1$，其一圈结果为
\begin{align} \label{eq:ps-c}
	&{\cal C}(\alpha,z^2\mu^2)\\
	=& \left[ 1+ {\alpha_sC_F\over 2\pi}\left({3\over2}\ln{z^2\mu^2e^{2\gamma_E}\over 4}+{3\over2}\right)\right]\delta(1-\alpha)\nn
	\\
	&+ {\alpha_sC_F\over 2\pi}\left\{\left({1+\alpha^2\over 1-\alpha}\right)^{[0,1]}_{+(1)}\left[-\ln{z^2\mu^2e^{2\gamma_E}\over 4}-1\right]\right.\nn\\
	&\left. - \left(4\ln(1-\alpha)\over 1-\alpha\right)^{[0,1]}_{+(1)}\!  +\!2(1+\sigma)(1-\alpha)\right\}\!\theta(\alpha)\theta(1-\alpha)
	\,,\nn
\end{align}
该结果也曾在~\cite{Ji:2017rah,Radyushkin:2017lvu,Radyushkin:2018cvn,Zhang:2018ggy,Li:2020xml}中计算与进一步研究。可以检验上述结果确实经由 Eq.~(\ref{eq:ftc-C}) 与动量空间的单圈匹配相联系。由于我们关心的是准 LF 关联与光线算符矩阵元 $O_{\gamma^+}(\lambda n)$ 的关系，Eq.~(\ref{eq:io-Q-fact}) 也可用~\cite{Balitsky:1987bk,Braun:1994jq,Braun:2007wv}的光线算符展开得到。

利用 OPE 或短距离展开，包含混合效应的胶子与单态夸克准 PDF 的精确因子化公式也已在坐标空间导出并被研究至单圈阶~\cite{Wang:2019tgg,Balitsky:2019krf}。

容易看出，LaMET 展开中的 $P^z\to\infty$ 极限与坐标空间因子化中的 $z\to 0$ 极限（保持 $\lambda=zP^z$ 有限）是等价的。但在实际格点 QCD 计算中，受限于具体设置的最大动量 $P^z_{\rm max}$，两种途径开始出现差别。

在 LaMET 的系统近似中，应当对所有可能的 $z$ 或 $\lambda$ 计算 $\tilde h (z, P^z_{\rm max})$，但实践中最大值 $\lambda_{\rm max} = z_{\rm max}P^z_{\rm max}$ 受格点体积以及大 $z$ 数据质量的限制。由于 QCD 禁闭，$\tilde h (z, P^z_{\rm max})$ 的关联长度约为 $\sim 1/\Lambda_{\rm QCD}$，在大 $z$ 呈指数衰减~\cite{Ji:2020brr}。因此，如果 $z_{\rm max}$ 足够大（例如质子尺寸 $\sim1$ fm），使 $\tilde h (z, P^z_{\rm max})$ 几乎衰减到零，那么截断的 Fourier 变换应快速收敛，截断效应主要影响小 $x \lesssim 1/\lambda_{\rm max}$ 区域的结果。若 $\tilde h (z, P^z_{\rm max})$ 呈指数衰减但在 $z_{\rm max}$ 处仍有明显非零值，则可以在 $z_{\rm max}$ 以外做物理动机的外推~\cite{Ji:2020brr} 再进行 Fourier 变换，这既消除了截断带来的非物理振荡，又只影响小 $x$ 区域。在动量空间，可以利用 LaMET 展开逐点计算 $x$ 处的 PDF，系统误差由 $\Lambda_{\rm QCD}^2/(xP^z_{\rm max})^2$ 与 $\Lambda_{\rm QCD}^2/((1-x)P^z_{\rm max})^2$ 控制，从而对某个 $x$ 区域 $[x_{\rm min}, x_{\rm max}]$ 给出达到目标误差的预言。

在坐标空间因子化中，把 $\tilde h(\lambda, z^2)$ 按 $z^2\Lambda^2_{\rm QCD}$ 展开；要使因子化公式有效，$z$ 必须保持在微扰区内。例如~\cite{Ji:2020brr} 的估计给出 $z_{\rm max}\sim 0.3$--$0.4$ fm。虽然有观察表明构造比值 $\tilde h (\lambda, z^2)$ 可能抵消 $z>0.4$ fm 处的高扭度贡献~\cite{Orginos:2017kos}，但这种抵消需要量化才能用于精密计算。在准 LF 关联范围有限的情况下，可以通过对 $x$ 依赖建模或更先进的技术（如贝叶斯分析~\cite{Bringewatt:2020ixn}或神经网络~\cite{Karpie:2019eiq,Cichy:2019ebf,DelDebbio:2020rgv}）提取 PDF，这与从实验数据提取 PDF 类似~\cite{Ma:2017pxb}，尽管拟合系统误差的量化颇具挑战。坐标空间因子化还可用于提取 PDF 的最低几阶矩~\cite{Karpie:2018zaz,Shugert:2020tgq,Joo:2020spy,Gao:2020ito}，其主要系统误差由 $z^2\Lambda_{\rm QCD}^2$ 控制。

迄今为止，关于准 PDF 与 pseudo-PDF 分析的比较研究非常有限~\cite{Bhat:2020ktg,Alexandrou:2020qtt}。两种策略的系统误差如何对比尚待观察。

%------------------------------------------------------------------------------
\subsection{非微扰重正化与匹配}
\label{sec:renorm-npr}
%------------------------------------------------------------------------------
准 PDF 非局域 Wilson 线算符的乘法可重正性使格点上的非微扰重正化成为可能，此后即可取连续极限。这是 LaMET 应用中的重要一步。一种做法是先对 Wilson 线作质量减除~\cite{Musch:2010ka,Ishikawa:2016znu,Chen:2016fxx,Zhang:2017bzy,Green:2017xeu,Green:2020xco}，再用格点微扰论或其他非微扰方案重正化剩余的 UV 发散。近年来更流行的是正规化无关动量减除（RI/MOM）方案~\cite{Constantinou:2017sej,Stewart:2017tvs,Alexandrou:2017huk,Chen:2017mzz,Liu:2018uuj}。在 $|z|\ll \Lambda_{\rm QCD}^{-1}$ 的坐标空间途径中，不同态之间准 LF 关联的比值~\cite{Radyushkin:2017cyf,Orginos:2017kos,Braun:2018brg,Li:2020xml}也被提议作为一种重正化方案。在大 $z$ 处，RI/MOM 方案与比值方案会在不同程度上引入额外的非微扰效应，可能扭曲原始准 LF 关联的 IR 性质。由于 Fourier 变换中大 $P^z$ 对长程贡献的压低，这种非微扰污染主要影响 $x$ 空间的端点区域；而在中等 $x$ 采用 RI/MOM 方案的现有 LaMET 计算（例如~\cite{Alexandrou:2018eet,Lin:2018qky}）受此类系统误差影响较小。不过，改用混合方案~\cite{Ji:2020brr}——在不同距离处利用不同方案的优势——即可避免上述复杂性。
下面依次讨论这些方案，重点放在混合重正化方案上。


在继续之前应注意：文献~\cite{Detmold:2005gg,Braun:2007wv,Ma:2017pxb}中的流—流关联不需要重正化或在格点上具有简单的重正化，尽管模拟它们的代价可能更高。此外，还有一种独特的方法：借助梯度流（gradient flow）用涂抹（smeared）费米场与规范场重新定义准 PDF~\cite{Monahan:2016bvm}。涂抹后的准 PDF 没有 UV 发散，在连续极限下保持有限，并可微扰匹配到 PDF~\cite{Monahan:2017hpu}。不过该方法尚待在格点上实现。

\subsubsection{Wilson 线质量减除方案}
\label{sec:renorm-npr-mass}

由于质量修正 $\delta m$ 包含全部线性 UV 发散，把它从准 PDF 中非微扰地减除是非常可取的。众所周知，Wilson 线重正化与静态夸克—反夸克势的附加重正化有关，即 $\delta m$，尤其在有限温度场论的语境中。对一个空间与时间方向尺寸为 $L\times T$ 的矩形 Wilson 圈，大 $T$ 时其真空期望值满足
\begin{align}
	\lim_{T\to\infty}W(L,T) = c(L) e^{-V(L)T}\,.
\end{align}
重正化静态势为
\begin{align}\label{eq:deltam1}
	V^{R}(L) = V(L) + 2\delta m\,,
\end{align}
$\delta m$ 可通过施加条件 $V^{R}(L_0)=0$（某个特定 $L_0$）来确定。或者，也可以从著名的弦势模型拟合 $\delta m$：
\begin{align}\label{eq:deltam2}
	V(L) = \sigma L- {\pi \over 12 L} -2\delta m\,.
\end{align}

除了用静态势确定 $\delta m$，还有人建议在辅助``重夸克''场论中用如下条件计算它~\cite{Green:2017xeu}：
\begin{align}\label{eq:deltam3}
	\delta m = {d\over dz} \ln{\rm Tr}\big\langle Q(x+zn_z)\bar{Q}(x)\big\rangle_{{\rm QCD}+Q} \Big|_{z=z_0}\,.
\end{align}

关于 $\delta m$ 的非微扰计算还有其他建议~\cite{Ji:2020brr}。例如，可以考察强子矩阵元、单夸克 Green 函数或固定规范下 $O_{\Gamma}(z,a)$ 真空期望值的渐近大 $z$ 行为。由所有这些矩阵元算出的 $\delta m$ 对格距 $a$ 有如下依赖：
\begin{equation}
	\delta m = m_{-1}(a)/a + m_0 \ ,
\end{equation}
其中 $m_{-1}$(a) 是幂次发散系数，与具体矩阵元无关；$m_0~\sim O(\Lambda_{\rm QCD})$ 有限且依赖外部态。$m_0$ 的确定可能相当不平凡；实际计算中可采用微调方法（如对 Wilson 费米子质量那样）寻找临界值，使最终结果在大 $P^z$ 极限收敛最快。

完成 Wilson 线质量减除后，$O_\Gamma(z,a)$ 中仍有对数 UV 发散。可以用格点微扰论把 $\delta m$ 减除后的 $O_\Gamma(z,a)$ 匹配到 $\MS$ 方案~\cite{Ishikawa:2016znu,Xiong:2017jtn,Constantinou:2017sej}，但其收敛性仍需在更高阶检验。文献~\cite{Green:2017xeu,Green:2020xco} 用类 RI/MOM 方案非微扰地重正化了这些对数发散。

Wilson 线质量减除已在格点上实现于~\cite{Musch:2010ka,Zhang:2017bzy,Green:2017xeu,Chen:2017gck,Alexandrou:2020qtt}。

\subsubsection{RI/MOM 方案}
\label{sec:renorm-npr-rimom}

RI/MOM 方案在格点 QCD 中已被广泛用于无幂次发散混合的局域复合夸克算符的重正化~\cite{Martinelli:1994ty}。它本质上是 QFT 中的动量减除方案，可以在格点上非微扰实现。对任意被乘法重正化的复合夸克双线性算符 $O^B=Z_OO^R$，RI/MOM 方案通过在其离壳夸克矩阵元上于减除标度 $\mu_{\tiny R}$ 施加如下条件来定义：
\begin{align} \label{eq:rimom}
	Z^{-1}_O \langle p|O^B|p\rangle \Big|_{p^2=-\mu_{\tiny R}^2}= \langle p|O|p\rangle_{\rm tree}\,.
\end{align}
下标``tree''指微扰论中的树图矩阵元。若 $\mu_{\tiny R} \gg \Lambda_{\rm QCD}$，\eq{rimom} 定义的 $Z_O$ 位于微扰区，可以按微扰论逐阶转换到 $\MS$ 方案。在这个意义上，$Z_O$ 并非字面意义的非微扰量，而是全阶可计算的量。

既然非局域夸克双线性算符 $O_\Gamma(z)$ 已被证明在坐标空间乘法可重正，也可以在 RI/MOM 方案中重正化它，然后把结果匹配到 $\MS$ 方案的 PDF~\cite{Constantinou:2017sej,Stewart:2017tvs}。在格点上，算符的离壳矩阵元由其带离壳夸克的截腿 Green 函数（即顶点函数）定义。对非局域 Wilson 线算符，后者是
\begin{align} \label{eq:Lambda}
	&\Lambda^{\Gamma}_0(z, a, p) \equiv
	\left[S^{-1}_0(p,a)\right]^\dagger \sum_{x,y}e^{ip\cdot (x-y)}
	\nn\\
	&\ \ \  \times\! \big\langle 0 \big|T\bigl[\psi_0(x,a) O^B_\Gamma(z,a) \bar{\psi}_0(y,a)\bigr]\big|0\big\rangle S^{-1}_0(p,a)\,,
\end{align}
其中 $S_0(p,a)$ 是裸夸克传播子，外部动量 $p$ 是格点上的欧氏动量。Green 函数不是规范不变量，需要固定规范（通常选 Landau 规范 $\partial\cdot  A=0$）；其规范依赖预期会被逐阶的匹配或方案转换所抵消。

计入夸克波函数重正化 $Z_q$（可在格点上独立确定~\cite{Martinelli:1994ty}）后，\eq{rimom} 修订为
\begin{align} \label{eq:rimomvertex}
	Z_q Z^{-1}_{O_\Gamma} \Lambda^{\Gamma}_0(z,a,p)\Big|_{p=p_{\tiny R}}= \Lambda^{\Gamma}_{\rm tree}(z,a,p) = \Gamma e^{ip_{\tiny R}\cdot z}\,.
\end{align}
由于 $O_\Gamma(z,a)$ 不是 $O(4)$ 协变的，RI/MOM 方案需用两个标度定义：一是 $\mu_{\tiny R}=|p_R|$，另一是 $p_{\tiny R}^z$。为简便计统记为 $p=p_{\tiny R}$。要在微扰区工作并控制 ${\cal O}\big(a^2\mu_{\tiny R}^2, a^2(p_{\tiny R}^z)^2\big)$ 阶的格点离散化效应，必须在窗口 $\Lambda_{\rm QCD}\ll \mu_{\tiny R}\ll a^{-1}$、$p_{\tiny R}^z \ll a^{-1}$ 内工作，这要求格距足够小。

由于夸克离壳，还可能出现与 EOM 算符的有限混合。结果是 \eq{rimomvertex} 一般不能作为矩阵方程满足，而需要一个投影算符 ${\cal P}$ 来定义离壳矩阵元：
\begin{align}
	\langle p| O^B_\Gamma |p\rangle = {\rm tr}\big[\Lambda^{\Gamma}_0(z,a,p) {\cal P}\big]\,,
\end{align}
以便计算重正化因子 $Z_{O_\Gamma}$。

然后，裸强子矩阵元 $\tilde{h}_B(z,P^z,a)$ 可以在坐标空间重正化为
\begin{align} \label{eq:rimomh}
	\tilde{h}_R(z,P^z, p_{\tiny R}^z,\mu_{\tiny R},a)=Z^{-1}_{O}(z,p_{\tiny R}^z,\mu_{\tiny R},a)\tilde{h}_B(z,P^z,a) \,,
\end{align}
在连续极限下，重正化矩阵元与 UV 调节子无关，因此在 RI/MOM 方案下用 DR 应得到同样结果：
\begin{align}\label{eq:continuum}
	&\tilde{h}_R(z,P^z, p_{\tiny R}^z,\mu_{\tiny R})=\lim_{a\to0}\tilde{h}_R(z,P^z, p_{\tiny R}^z,\mu_{\tiny R},a) \nn\\
	&\qquad = \lim_{\epsilon\to0}Z^{-1}_{O}(z,p_{\tiny R}^z,\mu_{\tiny R},\epsilon)\tilde{h}_B(z,P^z,\epsilon)\,,
\end{align}
这使我们能在连续微扰论中计算匹配系数。注意由于使用 DR，$Z_{O}$ 中不含 $\delta m$。

把上述重正化矩阵元 Fourier 变换到动量空间后，就可以求出准 PDF 的 RI/MOM 匹配系数~\cite{Stewart:2017tvs}。不同自旋结构的一圈匹配系数已在 \cite{Stewart:2017tvs,Liu:2018uuj,Liu:2018hxv} 得到，非极化情形的两圈结果见~\cite{Chen:2020ody}。另一种做法是先把 RI/MOM 矩阵元转换到 $\MS$ 或修改的 $\MS$ 方案~\cite{Constantinou:2017sej,Alexandrou:2019lfo}，然后再做 Fourier 变换与动量空间匹配。

\subsubsection{比值方案}
\label{sec:renorm-npr-ratio}

在坐标空间因子化中，$|z|\ll \Lambda^{-1}_{\rm QCD}$ 必须很小，而 $P^z$ 可以取任意值。此时文献~\cite{Radyushkin:2017cyf,Orginos:2017kos}的比值方案可以是格点重正化的有效选择。考虑比值
\beq \label{eq:ratio}
\tilde{h}(\lambda,z^2,a)/\tilde{h}(0,z^2,a)\,,
\eeq
其中分母是 $P^z=0$ 的非微扰矩阵元。由于 $\tilde{h}(\lambda,z^2,a)$ 与 $\tilde{h}(0,z^2,a)$ 来自同一格点系综、彼此相关，比值的统计误差可以降低。而且该比值无需在格点上进一步重正化，可直接取连续极限
\begin{align} \label{eq:ratio1}
	\lim_{a\to0}\frac{\tilde{h}(\lambda,z^2,a)}{\tilde{h}(0,z^2,a)} = \frac{\tilde{h}(\lambda,z^2)}{\tilde{h}(0,z^2)}\,,
\end{align}
文献~\cite{Radyushkin:2017cyf,Orginos:2017kos}称之为``约化 Ioffe 时间赝''分布。在 $\MS$ 方案中，$\tilde{h}(0,z^2\mu^2)$ 有小 $z$ 展开
\begin{align}
	\tilde{h}(0,z^2\mu^2) = C_0(z^2\mu^2) + {\cal O}(z^2M^2, z^2\Lambda_{\rm QCD}^2)\,,
\end{align}
其中同位旋矢量夸克 PDF 的最低矩 $a_0$ 平凡地为 1。若忽略全部幂次修正，则 $\tilde{h}(0,z^2\mu^2)$ 是微扰量，可视为一个重正化因子。因此 \eq{ratio1} 的比值仍满足与 \eqs{ope}{io-Q-fact} 类似的 OPE 或因子化公式，只需相应修改匹配系数~\cite{Radyushkin:2017lvu,Izubuchi:2018srq}：
\begin{align}
	{\cal C}^{\rm ratio}(\alpha,z^2\mu^2)= {\cal C}(\alpha,z^2\mu^2) - \delta(1-\alpha)C_0(z^2\mu^2)\,.
\end{align}

在比值方案的其他变体中，也有人建议用非局域 Wilson 线算符的真空矩阵元替代 $\tilde{h}(0,z^2,a)$~\cite{Braun:2018brg,Li:2020xml}，因为 UV 发散不依赖外部态。


\subsubsection{混合方案}
\label{sec:hybrid}

由于准 PDF 的因子化公式只在 $\MS$ 方案中被证明，对任何与 $\MS$ 非微扰不同的方案使用动量空间因子化并不合法。RI/MOM 与比值方案即属此类，因为把它们转换到 $\MS$ 的转换因子含 $z^2$ 的对数~\cite{Constantinou:2017sej,Izubuchi:2018srq}，当 $z\sim \Lambda_{\rm QCD}^{-1}$ 时需要把 $\alpha_s$ 跑到 IR 区。相比之下，带波函数重正化的 Wilson 线质量减除方案本质上等同于 $\MS$，不会引入额外的 IR 效应。

然而 Wilson 线质量减除方案也有缺点：在格点上，由于 $z\sim a$ 处的离散化效应，格点方案无法重现非局域算符 $\MS$ 矩阵元的短距离 $\ln z^2$ 行为；而这种离散化效应恰好在 RI/MOM 或比值方案中相消。为了兼取两类方案的优势，文献~\cite{Ji:2020brr}提出了混合方案，为在所有 $z$ 重正化准 LF 关联提供了一条可行路径。


具体而言，对 $|z|\le z_{\rm S}$（$z_{\rm S}$ 小于 OPE 中领头扭度近似失效的距离），把准 LF 关联重正化为
\begin{align}\label{eq:ratio2}
	\frac{\tilde h(z,a,P^z)}{Z_X(z,a)}\,,
\end{align}
其中``$X$''可为 RI/MOM 或比值方案。

对 $|z|>z_{\rm S}$，则应用 Wilson 线质量减除
\begin{align}
	{\tilde h(z,a,P^z)}e^{-\delta m |z|}Z_{\rm hybrid}\,,
\end{align}
其中 $Z_{\rm hybrid}$ 代表波函数与顶点重正化，可通过在 $z=z_{\rm S}$ 处施加连续性条件非微扰确定：
\begin{align}
	Z_{\rm hybrid} e^{-\delta m |z_{\rm S}|}{\tilde h(z,a,P^z)}  = \frac{{\tilde h(z,a,P^z)}}{Z_X(z_{\rm S},a) }\,,
\end{align}
从而
\begin{align}
	Z_{\rm hybrid}(z_{\rm S},a) = e^{\delta m |z_{\rm S}|}/{Z_X(z_{\rm S},a)  }\,.
\end{align}
如此一来只需计算 $\delta m$。注意最终结果应与 $z_{\rm S}$ 无关，因此应尝试多个数值，寻找使结果变化最平缓的最优点。

混合重正化准 PDF 的微扰匹配可以相应导出。以 $Z_X$ 取比值方案的零动量矩阵元为例，$O(\alpha_s)$ 匹配已在~\cite{Ji:2020brr} 导出：
\begin{align}\label{eq:hybridm2}
	&C_{\rm hybrid}(\xi, \mu^2/(p^z)^2, z_{\rm S}^2\mu^2) = C_{\rm ratio}(\xi, \mu^2/(p^z)^2) \nn\\
	&\quad+{\alpha_sC_F\over 2\pi}{3\over 2} \left[-{1\over |1-\xi|_+}+ \frac{2 \text{Si}((1-\xi) \lambda_{\rm S})}{\pi  (1-\xi)} \right]\,,
\end{align}
其中 $C_{\rm ratio}$ 见~\cite{Izubuchi:2018srq}，$\xi=y/x$，$\lambda_{\rm S}=z_{\rm S}p^z$，$p^z=xP^z$ 是部分子动量。加号函数定义为
\begin{align}
	{1\over |1-\xi|_+} &\equiv \lim_{\beta\to0^+}\left[ {\theta(|1-\xi|-\beta)\over |1-\xi|} + 2\delta(1-\xi)\ln\beta\right] \,.
\end{align}

由于格点体积有限以及大 $z$ 处信噪比恶化，可用的格点数据必须在 $z_{\rm L}$ 处截断。如 \sec{renorm-coord} 所讨论，准 LF 关联的关联长度为 $\xi_z\sim \Lambda_{\rm QCD}^{-1}$，在大 $z$（$\sim 1$ fm）呈指数衰减。若 $z_{\rm L}$ 不够大、准 LF 关联仍有相当的非零值，则在 $z_{\rm L}$ 截断的直接 Fourier 变换会导致非物理振荡及其他系统误差。

为改善这一点，混合方案建议做 $z\to\infty$ 外推~\cite{Ji:2020brr}。当 $P^z$ 不太大且格点矩阵元在 $z_{\rm L}$ 附近呈指数行为时，可用 $\sim e^{-z/\xi_z}$ 形式外推，也可叠加一些代数行为以更好地反映 $z$ 依赖。若 $P^z$ 很大，信噪比变差，$z_{\rm L}$ 更小；此时准 LF 关联尚未表现出指数衰减而被领头扭度贡献主导，可用代数衰减形式外推。鉴于当代计算资源可使 $\lambda_{\rm L} =z_{\rm L}P^z$ 达到相当大的值，外推只影响很小的 $x$ 区域——那本来就是 LaMET 展开控制欠佳的区域。

总之，混合方案为所有 $z$ 的准 LF 关联提供了恰当的重正化，使我们能够通过动量空间的 LaMET 展开对 $x\in[x_{\rm min}, x_{\rm max}]$ 受控地计算 PDF。因此，我们期待它在未来 LaMET 的 PDF 计算中扮演主导角色。

\subsection{总胶子螺旋度 $\Delta G$ 与横性 PDF}
\label{sec:deltaG}

除共线 PDF 外，LaMET 的首个应用是胶子螺旋度对质子自旋的贡献 $\Delta G$~\cite{Ji:2013fga}。在质子自旋的朴素求和规则~\cite{Jaffe:1989jz}中，$\Delta G$ 与一个非局域光锥关联算符的矩阵元相联系~\cite{Manohar:1990jx}：
\begin{align}\label{eq:sginv}
	\Delta G &\!=\! \langle PS\big|i\!\int\! \frac{dxd\lambda}{2\pi xP^+} e^{i\lambda x} \! F^{+\alpha}(0) W(0,\!\lambda n)\tilde F_{\alpha}^{~+} (\lambda n)\big|PS\rangle\,,
\end{align}
在光锥规范 $A^+=0$ 下它约化为
\begin{align}
	\Delta G &= \langle PS| \big(\vec{E}\times \vec{A}\big)^z|PS\rangle /(2P^+)\,.
\end{align}

在 LaMET 框架内，可以从一个静态``胶子自旋''算符出发：即在保持 IMF 极限下胶子场横向极化的时间无关规范中固定的 $\vec{E}\times \vec{A}$。例如库仑规范 $\vec{\nabla}\cdot \vec{A}=0$、轴向规范 $A^z=0$ 与 $A^0=0$ 都是可行选择~\cite{Hatta:2013gta}。

在库仑规范与 $\MS$ 方案中，一圈阶在壳有质量夸克态中的静态``胶子自旋'' $\Delta \widetilde G$ 为~\cite{Chen:2011gn,Ji:2013fga}
\begin{align}  \label{eq:coulomb}
	&\Delta \widetilde G(P^z, \mu)  (2S^z)=\left.\langle PS| \epsilon^{ij}_\perp F^{i0}A^j|PS\rangle_q \right\arrowvert_{\vec{\nabla}\cdot \vec{A}=0}\\
	&\qquad\qquad = \frac{\alpha_sC_F}{4\pi} \left[ {5\over 3}\ln{\mu^2\over m^2} -\frac{1}{9} + \frac{4}{3}\ln \frac{(2P^z)^2}{m^2}\right](2S^z) \,,\nn
\end{align}
其中下标 $q$ 表示夸克，$S^\mu$ 是自旋矢量。共线发散由有限的夸克质量 $m$ 调节。

如果循着~\cite{Weinberg:1966jm}的流程，在 UV 正规化之前先取 $P^z\to\infty$ 极限~\cite{Ji:2013fga}，则
\begin{align} \label{eq:coulombimf}
	\Delta \widetilde{G}(\infty, \mu) (2S^z)&= \left. \langle PS| \epsilon^{ij}_\perp F^{i0}A^j|PS\rangle_q\right\arrowvert_{\vec{\nabla}\cdot \vec{A}=0}\nn\\
	&=\frac{\alpha_sC_F}{4\pi}\left(3\ln{\mu^2\over m^2}+7\right) (2S^z)\,,
\end{align}
这与光锥胶子螺旋度 $\Delta G(\mu)$ 完全一致~\cite{Hoodbhoy:1998bt}。因此，尽管 UV 发散不同，$\Delta \widetilde G(P^z,\mu)$ 与 $\Delta G(\mu)$ 的共线发散完全相同，从而允许在二者之间进行微扰匹配。

联系 $\Delta\widetilde G(P^z,\mu)$ 与 $\Delta G$、$\Delta \Sigma$ 的完整因子化公式是
\begin{align} \label{eq:deltaGmatching}
	\Delta\widetilde G(P^z,\mu) &= Z_{gg}(P^z/\mu)\Delta G(\mu) \nn\\
	&\qquad\qquad+ Z_{gq}(P^z/\mu)\Delta\Sigma(\mu) +...\,,
\end{align}
其中 $...$ 是被 $1/P^z$ 压低的幂次修正；库仑规范的匹配系数 $Z_{gg}$ 与 $Z_{gg}$ 已在一圈计算~\cite{Ji:2014lra}。

此外，也可以按 \sec{renorm} 的因子化公式先计算胶子螺旋度 PDF $\Delta g(x)$，再对 $x$ 积分得到 $\Delta G$。\\

在领头扭度下，除了前面讨论的非极化与螺旋度 PDF，还存在定义为~\cite{Jaffe:1991kp,Jaffe:1991ra}
\begin{equation}
	h_1(x)\! =\! \frac{1}{2P^+}\!\int \frac{d\lambda}{2\pi}
	e^{-i\lambda x} \langle PS_\perp|\overline{\psi}(0)\gamma^+\gamma_\perp\gamma_5\psi(\lambda n)|PS_\perp\rangle \,.
\end{equation}
的横性（transversity）PDF。$h_1(x)$ 简单地数出横向极化质子中携带动量份额 $x$ 的横向极化夸克的数目。该分布的第一矩对应所谓张量荷 $\delta q$，它是手征奇算符的矩阵元。$h_1(x)$ 可通过 Drell-Yan 过程中的横—横自旋不对称~\cite{Ralston:1979ys,Jaffe:1991kp,Jaffe:1991ra}或 SIDIS 中的 Collins 单自旋不对称（横性 TMDPDF 与手征奇的 TMD 碎裂函数耦合）~\cite{Collins:1992kk}来测量。目前关于横性 PDF 的实验结果非常有限~\cite{Barone:2001sp,Kang:2015msa,Lin:2017stx,Radici:2018iag,Cammarota:2020qcw}，海夸克贡献尤其如此~\cite{Chang:2014jba}，因此这是格点 QCD 计算能够发挥重要作用的一个方向。

$h_1(x)$ 的 LaMET 计算直截了当，因为其非局域算符的重正化与非极化情形相同；其 $\MS$ 与 RI/MOM 方案的一圈匹配均已计算~\cite{Alexandrou:2018eet,Liu:2018hxv}。$h_1(x)$ 的首批格点计算见于~\cite{Chen:2016utp,Alexandrou:2018eet,Liu:2018hxv}，将在 \sec{lattice} 更详细讨论。


%------------------------------------------------------------------------------
